%------------------------------------------------------------------------------%

\usepackage{apresentacao}
\usepackage{mathconfig}
\usepackage{tabto}

\makeatletter
\graphicspath  {{./figs/}}
\def\input@path{{./figs/}}
\makeatother

%------------------------------------------------------------------------------%
\title {Funções e Derivadas}
\author{\large Luis Alberto D'Afonseca}
\date  {Integrais e Séries}

%------------------------------------------------------------------------------%
\begin{document}

%------------------------------------------------------------------------------%
\begin{frame}
    \titlepage
\end{frame}

%------------------------------------------------------------------------------%
\section{Função}

\begin{frame}{Função}
  \emph{Função} de $D$ em $C$
  \[
    \begin{array}{rccl}
      f\colon & D & \to & C \\
              & x & \to & y=f(x)
    \end{array}
  \]
  \pause
  associa a cada elemento, $x$, do \emph{domínio}, $D$, \\
  um único elemento, $y$, do \emph{contradomínio} $C$

  \vspace{\baselineskip}
  \pause
  \emph{Imagem} \\
  \pause
  Conjunto dos elementos $y$ associados a algum $x$
\end{frame}

%------------------------------------------------------------------------------%
\begin{frame}{Função Real}
  \emph{Função} de $\R$ em $\R$
  \[
  \begin{array}{rccl}
    f\colon & \R & \to & \R \\
    & x & \to & y=f(x)
  \end{array}
  \]
  domínio e contra domínio são $\R$
\end{frame}

%------------------------------------------------------------------------------%
\begin{frame}{Funções Comuns}
  \begin{description}[<+->]
    \item[Função constante]  \tabto{30mm} \(f(x) = c\)      \\[5mm]
    \item[Função afim]       \tabto{30mm} \(f(x) = ax + b\) \\[5mm]
    \item[Função potência]   \tabto{30mm} \(f(x) = x^n\)    \\[5mm]
    \item[Função polinômial] \tabto{30mm} \(f(x) = a_nx^n + a_{n-1}x^{n-1}
                                                 + \cdots + a_1 x + a_0\)
  \end{description}
\end{frame}

%------------------------------------------------------------------------------%
\begin{frame}{Funções Exponencial e Logarítmica}
  \begin{description}[<+->]
    \item[Função exponencial]          \tabto{30mm} \(f(x) = a^x\)       \\[5mm]
    \item[Função exponencial natural]  \tabto{30mm} \(f(x) = e^x\)       \\[5mm]
    \item[Função logarítmo]            \tabto{30mm} \(f(x) = \log_b(x)\) \\[5mm]
    \item[Função logarítmo natural]    \tabto{30mm} \(f(x) = \ln(x) = \log_e(x)\)
  \end{description}
  \onslide<+->{
  \begin{align*}
    y & = \log_b(x) & \quad\Rightarrow\qquad b^y = x & \hspace{90mm}~\\[3mm]
    y & = \ln(x)    & \quad\Rightarrow\qquad e^y = x &
  \end{align*}
  \vspace{-\baselineskip}
  }
\end{frame}

%------------------------------------------------------------------------------%
\begin{frame}{Funções Trigonométricas}
  \begin{description}[<+->]
    \item[Seno]       \tabto{30mm} \(f(x) = \sin(x)\)       \\[5mm]
    \item[Cosseno]    \tabto{30mm} \(f(x) = \cos(x)\)       \\[5mm]
    \item[Tangente]   \tabto{30mm} \(f(x) = \tan(x) = \frac{\sin(x)}{\cos(x)}\) \\[5mm]
    \item[Cotangente] \tabto{30mm} \(f(x) = \cot(x) = \frac{\cos(x)}{\sin(x)} = \frac{1}{\tan(x)}\) \\[5mm]
    \item[Secante]    \tabto{30mm} \(f(x) = \sec(x) = \frac{1}{\cos(x)} \)      \\[5mm]
    \item[Cossecante] \tabto{30mm} \(f(x) = \csc(x) = \frac{1}{\sin(x)} \)
  \end{description}
\end{frame}

%------------------------------------------------------------------------------%
\section{Derivadas}

%------------------------------------------------------------------------------%
\begin{frame}{Derivada de uma Função Real}
  A \emph{Derivada} de uma função real $f(x)$ em relação
  a variável $x$
  \pause
  é a função $f'$\\
  cujo valor em $x$ é
  \[
    f'(x) = \lim_{h\to 0} \frac{\,f(x+h)-f(x)\,}{h}
  \]
  desde que o limite exista.
\end{frame}

%------------------------------------------------------------------------------%
\begin{frame}{Interpretação Geométrica da Derivada}
  A \emph{derivada} de $f$ no ponto $x$ é \\[3mm]
  \pause
  o \emph{coeficiente angular}
  da reta tangente ao gráfico de $f$ no ponto $x$
\end{frame}

%------------------------------------------------------------------------------%
\begin{frame}{Função Derivada}
  A \emph{função derivada} é \\[3mm]
  \pause
  a função \(f'(x)\) que retorna o valor da derivada de $f$ em cada ponto $x$
\end{frame}

%------------------------------------------------------------------------------%
\begin{frame}{Algumas Funções Derivada}
\begin{columns}
\column{0.5\linewidth}
\begin{gather*}
  \onslide<+->{ \frac{dc}{dx} = 0 \\[3mm]}
  \onslide<+->{ \frac{dx}{dx} = 1 \\[3mm]}
  \onslide<+->{ \frac{d}{dx} x^n = n x^{n-1}}
\end{gather*}
\column{0.5\linewidth}
\begin{gather*}
  \onslide<+->{ \frac{d}{dx} e^x = e^x \\[3mm]}
  \onslide<+->{ \frac{d}{dx} \ln\abs{x} = \frac{1}{x}}
\end{gather*}
\end{columns}
\end{frame}

%------------------------------------------------------------------------------%
\begin{frame}{Algumas Funções Derivada}
  \begin{columns}
    \column{0.5\linewidth}
    \begin{gather*}
      \onslide<+->{ \frac{d}{dx} \sin(x) = \cos(x) \\[3mm]}
      \onslide<+->{ \frac{d}{dx} \cos(x) = -\sin(x) \\[3mm]}
      \onslide<+->{ \frac{d}{dx} \tan(x) =  \sec^2(x) \\[3mm]}
    \end{gather*}
    \column{0.5\linewidth}
    \begin{gather*}
      \onslide<+->{ \frac{d}{dx} \cot(x) = -\csc^2(x) \\[3mm]}
      \onslide<+->{ \frac{d}{dx} \sec(x) =  \sec(x) \tan(x) \\[3mm]}
      \onslide<+->{ \frac{d}{dx} \csc(x) = -\csc(x) \cot(x)}
    \end{gather*}
  \end{columns}
\end{frame}

%------------------------------------------------------------------------------%
\begin{frame}{Propriedades da Derivada -- Regas de Derivação}
  $f$ e $g$ funções deriváveis, $c$ e $r$ constantes
  \begin{columns}
    \column{0.5\linewidth}
    \begin{enumerate}[<+->]
      \item \( \big(cf \big)' = cf' \) \\[3mm]
      \item \( \big(f+g\big)' = f' + g' \)\\[3mm]
      \item \( \big(f-g\big)' = f' - g' \)\\[3mm]
      \item \( \big(fg \big)' = f'g + fg' \)
    \end{enumerate}
    \column{0.5\linewidth}
    \vspace{\baselineskip}
    \begin{enumerate}[<+->]
      \setcounter{enumi}{4}
      \item \( \left(\frac{f}{\,g\,}\right)' = \frac{f'g - fg'}{g^2} \)\\[7mm]
      \item Se \(h(x) = f\big(g(x)\big) \) então  \[ h'(x) = f'\big(g(x)\big)\,g'(x) \]
    \end{enumerate}
  \end{columns}
\end{frame}

%------------------------------------------------------------------------------%
\begin{frame}
  \tableofcontents
\end{frame}

%------------------------------------------------------------------------------%
\end{document}
%------------------------------------------------------------------------------%
